\chapter{Part III}
Part III (a lossy codec based on the DCT and a comparison with MPEG
codecs) was developed using Claude Code (Anthropic).

We defined two rules right at the beginning, which shaped the entire
interaction:

\begin{itemize}
    \item The AI writes the code (C++17, CMake); the decisions are ours.
    \item The work is divided into small steps, and at the end of each
    step we test it ourselves before moving on.
\end{itemize}

\section*{What It Was Used For}

\begin{itemize}
    \item \textbf{Planning:} Dividing the codec into steps that compile
    and can be tested one at a time.

    \item \textbf{Code:} WAV reading/writing, fast MDCT (FFT),
    quantizer, Mid/Side, context-based Rice coding, file format,
    rate control, measurement and testing scripts.

    \item \textbf{Verification:} All automated tests were written by
    the AI (tools/system\_test.sh: system tests with sanitizers and
    fuzzing); we confirmed the results using our Part I tool.

    \item \textbf{Documentation:} READMEs.
\end{itemize}

\section*{Problems: What the AI Proposed, How It Was Verified, and What Changed}

Cases in which the AI's first proposal was wrong or incomplete.

\begin{table}[H]
\centering
\small
\begin{tabularx}{\textwidth}{@{}p{0.17\textwidth}p{0.22\textwidth}p{0.29\textwidth}X@{}}
\toprule
\textbf{Problem} &
\textbf{What the AI Proposed} &
\textbf{How It Was Verified} &
\textbf{Result} \\
\midrule

Comparing v1 with v2 &
Bitrate with the same quantization step &
With the same step, the SNR differed by 0.5 dB (deadzone) &
Comparison was changed to use the same SNR (BD-rate, Bjøntegaard
method): $-32.7\,\%$ in bits \\

Use of Mid/Side in the README &
``Used in 70--99 \% of blocks, never in sample02'' &
Measured using the encoder statistics (-v) &
Wrong: 30--96 \%, and $\sim$3 \% in sample02. Corrected \\

RDO Quantization &
Started implementing it immediately after the deadzone &
We interrupted it: it had not been requested and the step was closed &
Left as future work, with the explicit indication ``not measured'' \\

Automated Test Criterion &
``SNR $>$ 30 dB'' with the default step &
Failed with our quieter WAVs, even though the codec was working
correctly &
Incorrect criterion: SNR depends on volume. Replaced with an error
limit derived from the step, MSE $\leq (1-r)^2 \cdot q^2/3$,
justified by measurements \\

Test Script Argument &
Accepted only a directory &
With a file, it said ``no WAVs'' and skipped 4 groups of tests
without warning; the AI tried to guess other causes before seeing
the output &
Changed to accept either a file or a directory, and to stop with
a clear message in other cases \\

Alignment in the Comparison with MPEG &
Add an offset search to wav\_snr (MP3/AAC/Opus introduce a delay
at the beginning) &
We tested this first: with .mp3, .m4a and .opus, ffmpeg already
removes the delay (offset 0); only raw .aac was offset by
1024 samples &
Simpler solution, without changing src/: use these formats and
have the script verify alignment, stopping if it is not 0 \\

\bottomrule
\end{tabularx}
\end{table}

\section*{How We Verified the Results}

No result was accepted simply because the AI said so:

\begin{itemize}
    \item \textbf{Independent Tool:} The SNR measured by wav\_snr
    in the lossy codec was compared with our wav\_cmp from Part I,
    and the values match.

    \item \textbf{Lossless Limit:} With the minimum step, 8-bit and
    16-bit files are reconstructed exactly (maximum error 0).

    \item \textbf{System Tests:} All automated tests were written by
    the AI. There are 82 end-to-end tests, compiled with
    AddressSanitizer and UndefinedBehaviorSanitizer, including
    invalid inputs and 200 randomly corrupted compressed files.
    We also ran them with our own WAVs, which is how we discovered
    the incorrect test criterion.

    \item \textbf{Reproducibility:} All tables in Report 1 are
    generated by a script. Running the scripts twice produces the
    same numbers; only the timings vary, by about 10 \%.

    \item \textbf{Listening:} A script generates the encoded files
    and the codec noise at several bitrates, so that each decision
    can be heard and not just measured.
\end{itemize}

\section*{Critical Evaluation}

\subsection*{What Worked Well}

It was most useful where the work was extensive but well-defined:
writing the codec and the measurement scripts.

It suggested methods we were not familiar with: BD-rate (used in
video coding), ablation to measure each improvement, and sanitizers
and fuzzing for testing.

\subsection*{What Did Not Work So Well}

\begin{itemize}
    \item \textbf{Confident Claims Before Measuring.}
    Several statements (use of Mid/Side, size of the file format,
    duration of the files) were presented as facts and were incorrect.
    They were only corrected because we required a measurement for
    every number.

    \item \textbf{Doing More Than Requested.}
    It started implementing features we had not requested (the RDO
    quantizer, extra rules in the tests). We had to explicitly ask it
    to explain each change before applying it.

    \item \textbf{Guessing.}
    When asked about a failure it had not seen, it gave plausible but
    incorrect causes. The actual cause became clear as soon as we
    showed it the output.

\end{itemize}